\documentclass[11pt,a4paper]{report}
\usepackage[T1]{fontenc}
\usepackage[utf8]{inputenc}
\usepackage[margin=28mm]{geometry}
\usepackage{amsmath}
\usepackage{booktabs}
\usepackage{longtable}
\usepackage{array}
\usepackage{fancyhdr}
\usepackage[hidelinks]{hyperref}
\usepackage{xcolor}
\usepackage{tcolorbox}

\newcommand{\bytes}[1]{\texttt{#1}}

\newtcolorbox{note}[1]{colback=blue!4!white,colframe=blue!40!black,title={#1}}
\newtcolorbox{caution}[1]{colback=red!4!white,colframe=red!50!black,title={#1}}

\pagestyle{fancy}
\fancyhf{}
\fancyhead[L]{\small Technics SX-WSA1R \textendash{} System Exclusive Reference}
\fancyhead[R]{\small\leftmark}
\fancyfoot[C]{\thepage}
\renewcommand{\headrulewidth}{0.4pt}

\begin{document}

\begin{titlepage}
\centering
\vspace*{3cm}
{\Huge\bfseries Technics SX-WSA1R\par}
\vspace{6mm}
{\LARGE System Exclusive Reference\par}
\vspace{18mm}
{\large Acoustic Modeling Synthesizer Module\par}
\vspace{4mm}
{\large Applies equally to the SX-WSA1 keyboard\par}
\vspace{2mm}
{\large ROM set \textsc{wsa-a}~822 / \textsc{wsa-c}~230 / \textsc{wsa-d}~54\par}
\vspace{20mm}
{\Large\bfseries Unofficial\par}
\vspace{6mm}
\begin{minipage}{0.78\textwidth}
\centering
\normalsize
This document is not published, endorsed or checked by Technics or Panasonic.
It was written by inspecting a disassembly of the program ROMs --- the same
ROMs ship in the SX-WSA1 keyboard and the SX-WSA1R rack --- and checking the
result against a bulk dump captured from a real instrument. The ROM set is the
one the instrument's own \textsc{rom version} screen reports as the three
numbers above.
\end{minipage}

\vspace{12mm}
{\large\bfseries Draft of version 1 --- for evaluation by technical proof-readers\par}
\vspace{4mm}
{\large 21 September 2026\par}

\vspace{14mm}
\begin{minipage}{0.78\textwidth}
\centering
\normalsize
\textbf{How this was made.} The research and the writing were carried out by the
artificial-intelligence model \textbf{Claude Opus~5} (Anthropic), in its 1M-context configuration, working from the ROM
images, a captured bulk dump and the published Technics manuals. The work was directed and
coordinated by \textbf{Felipe Correa da Silva Sanches}
$\langle$\texttt{juca@members.fsf.org}$\rangle$.

Every measurement in this document is produced by a script committed beside it, so any
figure here can be re-derived rather than taken on trust. Readers are asked to treat the
whole of it as a draft under review: corrections are wanted, and the chapter on what this
reference does not cover is part of what needs checking.
\end{minipage}
\vfill
\end{titlepage}

\tableofcontents

\chapter{Overview}

The Technics SX-WSA1R exchanges data over MIDI System Exclusive. This reference describes
the messages it sends and accepts: their byte layout, the bulk-dump transfer protocol, the
universal messages it honours, and what it reports when a transfer fails.

\section{Applicability}
This reference describes operating system version \textbf{2.0}. An earlier version~1
exists; its behaviour may differ and is not covered.

\section{Ports}
The instrument has two sets of MIDI sockets, MIDI~1 and MIDI~2. Ordinary performance data
is handled on both. \textbf{Bulk dump transfers use MIDI~1 only} --- a dump offered on
MIDI~2 is not received, and the instrument sends dumps only from MIDI~1.

\section{Conventions}
Byte values are hexadecimal. A field written \bytes{\textit{nn}} takes one of the values
listed for it. Ellipses inside a message stand for a body whose length depends on the
message.

Every message begins with \bytes{F0} and ends with \bytes{F7}, as required by MIDI.

\begin{note}{Real-time messages are safe inside a transfer}
System Real Time bytes (\bytes{F8}--\bytes{FF}) may be interleaved anywhere inside a
System Exclusive message without disturbing it. Conversely, while the instrument is
receiving a System Exclusive message it ignores \bytes{F8} \textsc{timing clock},
\bytes{FA} \textsc{start}, \bytes{FB} \textsc{continue} and \bytes{FC} \textsc{stop} for
transport purposes.
\end{note}

\chapter{At a glance}\label{ch:glance}

Everything in this chapter is stated again, with its conditions, in the chapters that
follow. This page is for someone who already knows the protocol and wants the bytes.

\section{Frame}
\begin{center}
\bytes{F0} \quad \bytes{50} \quad \textsc{cmd} \quad \bytes{04} \quad \textit{md} \quad
\bytes{11} \quad [\,\textsc{adr}$\times$3 \quad \textsc{siz}$\times$3 \quad \textsc{dt}\ldots \quad
\textit{cn} \quad \textit{sm}\,] \quad \bytes{F7}
\end{center}

\bytes{04} is \textsc{pc}, the product; \textit{md} is \bytes{00} for the SX-WSA1 and
\bytes{01} for the SX-WSA1R, and both are accepted by both; \bytes{11} is \textsc{ver}, the
exclusive specification's version. \textsc{adr} and \textsc{siz} are 21-bit values in three
seven-bit bytes, most significant first. There is no device or unit number anywhere.

\section{Messages}
\begin{center}
\begin{tabular}{lll}
\toprule
\bytes{F0 50 21 04 \textit{md} 11 F7} & \textsc{hrq} & hand shake request \\
\bytes{F0 50 22 04 \textit{md} 11 F7} & \textsc{hrt} & hand shake routine \\
\bytes{F0 50 23 7E F7} & \textsc{ack} & acknowledge \\
\bytes{F0 50 24 7E F7} & \textsc{nak} & not acknowledged \\
\bytes{F0 50 25 \textit{lo} \textit{hi} F7} & \textsc{tmp} & tempo, $\mathit{lo}+16\,\mathit{hi}$; no checksum \\
\bytes{F0 50 27 7E F7} & \textsc{eok} & end of category \\
\bytes{F0 50 28 7E F7} & \textsc{end} & end of dump \\
\bytes{F0 50 29 7E F7} & \textsc{err} & abort \\
\bytes{F0 50 2A 7E F7} & \textsc{ful} & destination full \\
\bytes{F0 50 2B \ldots} & \textsc{drq} & request: \textsc{adr}, \textsc{siz}, no \textsc{dt} \\
\bytes{F0 50 2C \ldots} & \textsc{itr} & write one parameter \\
\bytes{F0 50 2D \ldots} & \textsc{btr} & bulk data, first frame \\
\bytes{F0 50 7E \ldots} & \textsc{cdd} & bulk data, every frame after \\
\midrule
\bytes{F0 7E 7F 09 01 F7} & & General MIDI System On \\
\bytes{F0 7E 7F 09 02 F7} & & General MIDI System Off \\
\bottomrule
\end{tabular}
\end{center}

\section{Rules}
\begin{description}
\item[Payload.] Each source byte goes out as two bytes, \textbf{high nibble first}:
  \bytes{F1} becomes \bytes{0F 01}. A payload is always an even number of bytes.
\item[Checksum.] $\mathit{sm} = (0 - S) \mathbin{\&} \mathtt{7F}$, where $S$ sums every
  byte from \bytes{50} through \textit{cn}. Equivalently everything after \bytes{F0},
  checksum included, sums to zero modulo 128. Verified on reception for
  \bytes{7E}, \bytes{2B}, \bytes{2C} and \bytes{2D} only.
\item[Continue.] \textit{cn} is \bytes{01} while more frames follow, \bytes{00} on the last.
\item[Size.] At most 256 bytes per message including \bytes{F0} and \bytes{F7}: 120 source
  bytes in a \textsc{btr} frame, 125 in a \textsc{cdd} frame.
\end{description}

\section{Bulk dump}
\begin{center}
\begin{tabular}{lll}
\toprule
\textsc{adr} & area & menu name \\
\midrule
\bytes{40 00 00} & \textsc{panel} & \textsc{system, part \& midi} \\
\bytes{20 00 00} & \textsc{sound memory} & \textsc{sound} \\
\bytes{50 00 00} & \textsc{combination} & \textsc{combination} \\
\bytes{60 00 00} & \textsc{sequencer} & \textsc{sequencer} (keyboard only) \\
\bottomrule
\end{tabular}
\end{center}

Send \textsc{drq} with one of those to ask for a category. The instrument must be showing
the \textsc{sysex bulk dump} screen to receive a dump, and bulk traffic works on the
MIDI~1 sockets only.

\section{Errors}
\textsc{error 40} wrong product, model or version. \textsc{error 41} anything wrong with
reception, including a missing acknowledgement. \textsc{error 42} an internal deadline, not
a fault at the far end. \textsc{error 21} the destination is full.

\chapter{Message framing}

\section{Identifiers}
The instrument accepts exactly two System Exclusive identifiers:

\begin{center}
\begin{tabular}{ll}
\toprule
identifier & meaning \\
\midrule
\bytes{50} & Matsushita (Technics) \\
\bytes{7E} & Universal Non-Real Time \\
\bottomrule
\end{tabular}
\end{center}

A message carrying any other identifier is tracked to its \bytes{F7} and discarded. No
error is reported for it.

\begin{caution}{There is no device or unit number after the identifier}
Many instruments place a device id immediately after the manufacturer identifier. The
SX-WSA1R does not. The byte following \bytes{50} is already the message-family byte
described in chapter~\ref{ch:messageset}. A unit field does appear later in some families,
but it is not where the convention would put it.
\end{caution}

\section{Termination}
A message ends at \bytes{F7}. If any other status byte in the range
\bytes{80}--\bytes{F6} arrives while a message is open, the message is abandoned and its
contents discarded.

\section{Length limit}
A message may not exceed 254 bytes. A longer message is rejected and, inside a bulk-dump
session, reported as a reception error.

\chapter{The Technics message set}\label{ch:messageset}

\section{Message families}
After \bytes{F0 50}, the third byte selects a message family. The instrument accepts
fourteen values:

\begin{center}
\bytes{21 22 23 24 25 27 28 29 2A 2B 2C 2D 7E 7F}
\end{center}

\bytes{26} is not accepted. The SX-WSA1R's own product identifier is \bytes{2D}; its
outgoing bulk-dump data messages all begin \bytes{F0 50 2D}.

\section{The model triple}
The families \bytes{21}, \bytes{22}, \bytes{2B}, \bytes{2C} and \bytes{2D} continue with
three fixed bytes:

\begin{center}
\bytes{04} \quad \bytes{\textit{nn}} \quad \bytes{11}
\end{center}

where \bytes{\textit{nn}} is \bytes{00} or \bytes{01}. A message whose triple the
instrument does not recognise is rejected with \textsc{error 40}.

\begin{caution}{The middle byte names the model, and only on transmission}
The keyboard transmits \bytes{04 00 11}; the rack transmits \bytes{04 01 11}. The rack
substitutes the middle byte as the message leaves, in every \bytes{21}, \bytes{22},
\bytes{2C} and \bytes{2D} message it sends, and recomputes the checksum of those that
carry one, so the substitution is invisible in the checksum. A bulk dump captured from a
real rack carries \bytes{04 01 11} throughout.

Reception is model-blind. Both values are accepted by both models, on every family that
carries the triple, and reach the same handler; a message is never refused, ignored or
treated differently because of the middle byte. A receiver must therefore accept both,
and a file dumped from one model loads into the other.
\end{caution}

\section{Short messages}
These messages are complete as shown. They form the control vocabulary of a bulk-dump
session.

\begin{center}
\begin{tabular}{ll}
\toprule
message & meaning \\
\midrule
\bytes{F0 50 23 7E F7} & acknowledge \\
\bytes{F0 50 24 7E F7} & not acknowledged \\
\bytes{F0 50 27 7E F7} & end of category \\
\bytes{F0 50 28 7E F7} & end of dump \\
\bytes{F0 50 29 7E F7} & abort \\
\bytes{F0 50 2A 7E F7} & destination full \\
\bytes{F0 50 21 04 \textit{nn} 11 F7} & enquiry \\
\bytes{F0 50 22 04 \textit{nn} 11 F7} & enquiry reply \\
\bottomrule
\end{tabular}
\end{center}

\section{Tempo}
The family \bytes{25} carries a tempo setting, in the range 40 to 300 beats per minute.
A value outside that range is treated as 120.

See chapter~\ref{ch:models} --- the tempo message is one of two things the rack does not
support.

\section{Parameter families}
The families \bytes{2B} and \bytes{2C} carry individual parameter traffic rather than bulk
transfer. Each addresses a large parameter space, and neither is enumerated in this
edition; see chapter~\ref{ch:limits}.

\chapter{Bulk dump}

Bulk dump transfers the instrument's internal data as System Exclusive. The instrument is
silent for the duration, and the transfer uses the MIDI~1 sockets only.

\section{Categories and what each transmits}
The \textsc{sysex bulk dump} screen offers five categories. Each sends one or more data
messages at fixed dump addresses.

\begin{center}
\begin{tabular}{llllr}
\toprule
category & address bytes & length bytes & length & blocks \\
\midrule
\textsc{system, part \& midi} & \bytes{40 00 00} & \bytes{00 00 20} & 32 & \\
                              & \bytes{40 00 20} & \bytes{00 12 60} & 2\,400 & \\
\addlinespace
\textsc{sound}                & \bytes{20 00 00} & \bytes{10 00 00} & 262\,144 & 32 \\
\addlinespace
\textsc{combination}          & \bytes{50 00 00} & \bytes{00 06 00} & 768 & \\
                              & \bytes{50 06 00} & \bytes{05 40 00} & 90\,112 & 16 \\
\addlinespace
\textsc{sequencer}            & \bytes{60 00 00} & \bytes{00 18 00} & 3\,072 & \\
                              & \bytes{60 18 00} & \bytes{01 70 00} & 30\,720 & \\
                              & \bytes{62 08 00} & \multicolumn{2}{l}{set at send time} & \\
\addlinespace
\textsc{total keyboard}       & \multicolumn{4}{l}{the four categories in the order above} \\
\bottomrule
\end{tabular}
\end{center}

Lengths are in source bytes; on the wire each occupies twice as many, because of the
payload encoding described below.

The \emph{blocks} column is internal structure, not extra messages: \textsc{sound} sends
one header covering all 262\,144 bytes, which is 32 repeats of an 8\,192-byte layout, and
\textsc{combination}'s second block one header covering 90\,112 bytes as 16 repeats of
5\,632. Those are the instrument's 32 sounds and 16 combinations.

\textsc{total keyboard} runs the four categories one after another in the order listed,
each closed by its own end-of-category message, with a single end-of-dump message after
the last.

The \textsc{sequencer} category is the only one whose final message has no fixed length ---
it carries as much recorded material as the instrument holds.

\begin{note}{The dump address is not a memory address}
Within one category the addresses are consecutive: a block's address is the previous
block's address plus its length. Across categories they are not. \textsc{sound} and
\textsc{combination} are held side by side inside the instrument, yet their dump addresses
are 786\,432 apart; other pairs are the reverse.

Treat the address and length bytes together as a fixed name for one of the eight blocks.
The instrument does the same --- on reception it compares them against a fixed list and
performs no arithmetic on them.
\end{note}

\section{Data message}
\begin{center}
\bytes{F0 50 2D 04 \textit{nn} 11}\quad
\bytes{\textit{a2 a1 a0}}\quad
\bytes{\textit{l2 l1 l0}}\quad
\bytes{\textit{data\ldots}}\quad
\bytes{\textit{c}}\quad
\bytes{\textit{sum}}\quad
\bytes{F7}
\end{center}

\begin{description}
\item[\bytes{\textit{nn}}] unit field, \bytes{00} or \bytes{01}.
\item[\bytes{\textit{a2 a1 a0}}] destination address, three bytes, most significant first,
  seven bits each: the value is $a_2 \times 128^2 + a_1 \times 128 + a_0$.
\item[\bytes{\textit{l2 l1 l0}}] length in \emph{source} bytes, encoded the same way.
\item[\bytes{\textit{data}}] the payload, encoded as described below.
\item[\bytes{\textit{c}}] continuation flag: \bytes{01} if further messages follow,
  \bytes{00} on the last.
\item[\bytes{\textit{sum}}] checksum.
\end{description}

A message is filled to at most 252 bytes before being sent; longer transfers are split.

\section{Payload encoding}
Each source byte is sent as two message bytes, \textbf{high nibble first}:
\[ b \;\longrightarrow\; (b \gg 4),\quad (b \mathbin{\&} \mathtt{0F}) \]
So source byte \bytes{F1} appears on the wire as \bytes{0F 01}, and the payload occupies
twice as many bytes as the length field states. A receiver rejoins them as
$((\text{first} \mathbin{\&} \mathtt{0F}) \ll 4) \mathbin{|} (\text{second} \mathbin{\&}
\mathtt{0F})$.

An odd number of payload bytes is invalid.

\section{Checksum}
Let $S$ be the eight-bit sum of every byte from the identifier \bytes{50} through the
continuation flag inclusive. Then
\[ \mathit{sum} \;=\; (0 - S) \mathbin{\&} \mathtt{7F} \]
The leading \bytes{F0}, the checksum byte itself and the trailing \bytes{F7} take no part
in the sum. Equivalently, the sum of every byte after \bytes{F0} including the checksum is
zero modulo 128.

\begin{caution}{The checksum is not verified for every family}
On reception the checksum is checked only for families \bytes{7E}, \bytes{2B},
\bytes{2C} and \bytes{2D}. For the other accepted families a wrong checksum passes
unremarked.
\end{caution}

\section{Requesting a dump}
A dump can be asked for over MIDI; the front-panel \textsc{send} button is not the only
way to start one. The request names a category by its region byte:

\begin{center}
\bytes{F0 50 2B 04 \textit{nn} 11 \textit{area} 00 00 \textit{x x x} F7}
\end{center}

\begin{center}
\begin{tabular}{ll}
\toprule
\bytes{\textit{area}} & category requested \\
\midrule
\bytes{40} & \textsc{system, part \& midi} \\
\bytes{20} & \textsc{sound} \\
\bytes{50} & \textsc{combination} \\
\bytes{60} & \textsc{sequencer} \\
\bottomrule
\end{tabular}
\end{center}

The three bytes in the length position are ignored; any values are accepted. The
instrument then runs exactly the transfer the corresponding menu row would have run.

There is no request for \textsc{total keyboard}: the four categories must be asked for one
at a time.

\section{Opening a transfer}
A dump begins with a two-step exchange. Each step is attempted up to three times.

\begin{enumerate}
\item The sender transmits \bytes{F0 50 21 04 00 11 F7} and expects
  \bytes{F0 50 22 04 \textit{nn} 11 F7}.
\item The sender transmits \bytes{F0 50 22 04 00 11 F7} and expects
  \bytes{F0 50 23 7E F7}.
\end{enumerate}

If the exchange succeeds, every subsequent data message must be acknowledged. If all three
attempts of the first step fail, \textbf{the dump proceeds anyway, without
acknowledgements}, at a fixed pace of roughly 51\,ms between messages.

\section{Acknowledgement and retry}
With acknowledgements in force, each data message must be answered with
\bytes{F0 50 23 7E F7}. Any other reply causes the same message to be retransmitted, up to
three attempts, after which the transfer is abandoned.

A reply is awaited for about 5.1 seconds; during the opening exchange the wait is about
2.05 seconds.

\section{Continuation}
Data messages after the first open with \bytes{F0 50 7E}, followed by payload,
continuation flag, checksum and \bytes{F7}.

\section{Closing}
At the end of each category the sender transmits \bytes{F0 50 27 7E F7} and waits for an
acknowledgement. At the end of the whole transfer it transmits \bytes{F0 50 28 7E F7}. If
the transfer was abandoned after three failed retries it transmits \bytes{F0 50 29 7E F7}.

\section{Receiving}
When the SX-WSA1R is the receiver it answers each data message with
\bytes{F0 50 23 7E F7} to accept, \bytes{F0 50 24 7E F7} to reject, and
\bytes{F0 50 2A 7E F7} when the incoming data is larger than the destination.

While a bulk-dump session is open the instrument accepts only the \bytes{50} identifier;
universal messages are not honoured until the session ends.

\begin{caution}{The instrument must be on the SYSEX BULK DUMP screen to receive}
A dump cannot be pushed to an idle instrument. Reception happens only while that screen is
displayed; a dump sent at any other time is parsed, discarded, and not reported.

A dump offered on the MIDI~2 input is never received either, at any time. This is what lies
behind the published instruction to use MIDI~1 for bulk transfers.
\end{caution}

Seven of the eight patterns are pinned on all six bytes, address and length together.
Sending a correct address with a shortened length is not a short transfer --- it is an
unrecognised message, refused with \textsc{error 41}. Only \textsc{sequencer} block~3
has its length read as a number, because that block's size varies. An arbitrary
destination cannot be written.

The one length the message itself controls is \textsc{sequencer} part~3, which is capped at
330\,752 bytes; beyond that, and whenever the destination fills mid-message, the
instrument answers \bytes{F0 50 2A 7E F7} and reports \textsc{error 21}.

\begin{caution}{Memory-full is treated as end-of-dump}
When the SX-WSA1R is \emph{sending} and the receiver answers \bytes{F0 50 2A 7E F7} to
report that its memory is full, the instrument treats that exactly as it treats
\bytes{F0 50 28 7E F7}, end of dump. It stops, and it reports success. A sender that needs
to know whether the data arrived cannot rely on the instrument to notice.
\end{caution}

\chapter{What a block contains}\label{ch:blocks}

This reference describes the protocol that carries a dump. The meaning of the data inside a
block is outside its scope, with one exception worth stating: every block opens with a
fixed signature, and two of them carry readable names. Together these let a program check
a stored dump before sending it back, and list what is in it, without understanding the
rest.

\section{Signatures}\label{sec:signatures}
The bytes below are the first bytes of the reassembled block --- after the payload has been
rejoined from its nibbles, not as they appear on the wire. They were read out of a dump
from a real instrument, and the two rows of each entry are the same sixteen bytes shown
twice, as hexadecimal and as text.

%% GENERATED by notes/sysex-probes/sysex_block_signatures.py --tex
%% Regenerate after any change to the capture or the reassembler

\begin{center}
{\small
\begin{tabular}{llll}
\toprule
block & \textsc{adr} & first sixteen bytes & as text \\
\midrule
\textsc{system, part \& midi} 1 & \bytes{40 00 00} & \bytes{5A 5A 01 00 57 41 30 00} & \texttt{ZZ..WA0.} \\
 & & \bytes{00 00 A0 00 00 00 00 00} & \texttt{........} \\
\textsc{system, part \& midi} 2 & \bytes{40 00 20} & \bytes{78 10 20 20 20 20 20 20} & \texttt{x.~~~~~~} \\
 & & \bytes{20 20 20 20 20 20 20 20} & \texttt{~~~~~~~~} \\
\textsc{sound} & \bytes{20 00 00} & \bytes{57 53 41 20 53 4F 55 4E} & \texttt{WSA~SOUN} \\
 & & \bytes{44 20 52 41 4D 20 53 30} & \texttt{D~RAM~S0} \\
\textsc{combination} 1 & \bytes{50 00 00} & \bytes{5A 5A 5A 5A 00 00 57 53} & \texttt{ZZZZ..WS} \\
 & & \bytes{41 31 20 20 01 00 00 02} & \texttt{A1~~....} \\
\textsc{combination} 2 & \bytes{50 06 00} & \bytes{78 10 20 53 70 61 63 65} & \texttt{x.~Space} \\
 & & \bytes{20 49 6E 66 69 6E 69 74} & \texttt{~Infinit} \\
\bottomrule
\end{tabular}
}
\end{center}


A file whose first bytes do not match the block it claims to be is not a dump of that
category.

\section{Names}
The \textsc{sound} block continues, immediately after its sixteen-byte signature, with
sixteen-byte names in the instrument's own character set, padded with spaces. The second
\textsc{combination} block and the second \textsc{system, part \& midi} block use a
different framing: a two-byte prefix, then a sixteen-byte name, then two further bytes.

A librarian can therefore show the user what a stored file holds --- ``\texttt{SWEEP PAD}'',
``\texttt{Space Infinity}'' --- without decoding anything else.

\begin{note}{Reassembly is self-checking}
A block reassembled correctly comes out at exactly the length its own header declared. A
wrong nibble order, or a payload taken with its trailing flag and checksum still attached,
does not. Checking the length is a cheap way to confirm a receiver is unpacking correctly
before trusting anything it produces.
\end{note}

\chapter{Parameter messages}\label{ch:params}

Two message families address individual parameters rather than transferring blocks.
Family \bytes{2C} writes; family \bytes{2B} requests. Both share one address space.

\begin{center}
\bytes{F0 50 2B 04 \textit{nn} 11 \textit{r} \textit{s} \textit{p} \ldots F7}
\end{center}

The three bytes after the model triple select what is addressed.

\section{Parameter area}
When \bytes{\textit{r}} is \bytes{00}, \bytes{\textit{s}} selects a block and
\bytes{\textit{p}} selects a parameter within it.

\begin{center}
\begin{tabular}{llr}
\toprule
\bytes{\textit{s}} & block & parameters \\
\midrule
\bytes{00} & common block 0 & 16 \\
\bytes{01} & common block 1 & 6 \\
\bytes{08} & write only & 4 \\
\bytes{10} & common block 10 & 20 \\
\bytes{11} & common block 11 & 4, plus a further 68 values \\
\bytes{20}--\bytes{3F} & part 0 to part 31 & 57 each \\
\bytes{60} & single parameter & 1 \\
\bottomrule
\end{tabular}
\end{center}

The part blocks are regular: \bytes{\textit{s}} is \bytes{20} plus the part number, and all
32 parts take the same 57 parameters. A few parameters carry more than one data byte.

Under \bytes{\textit{s}} = \bytes{11}, values of \bytes{\textit{p}} from \bytes{03} upward
are admitted in three runs --- \bytes{20}--\bytes{36}, \bytes{40}--\bytes{56} and
\bytes{60}--\bytes{75}. Any other value in that block is accepted by the parser and then
discarded, producing no error and no effect.

\section{Whole-area request}
Family \bytes{2B} also requests an entire category, as described in chapter~4. Here
\bytes{\textit{r}} takes the region byte rather than \bytes{00}.

\section{A third region}
Both families accept \bytes{\textit{r}} values of \bytes{10}, \bytes{18} and \bytes{19},
which carry an address and a count. What they address is not covered by this edition.

\chapter{The sound areas}\label{ch:sound}

\textsc{normal sound} and \textsc{drum sound} are the two areas of the address map that
hold an instrument sound rather than a setting. They are reached the same way as any other
parameter, with \textsc{itr} to write and \textsc{drq} to request.

\begin{center}
\begin{tabular}{lll}
\toprule
\textsc{adr} & address & area \\
\midrule
\bytes{10 00 00} & \bytes{040000} & \textsc{normal sound}, individual \\
\bytes{18 00 00} & \bytes{060000} & \textsc{drum sound}, individual \\
\bottomrule
\end{tabular}
\end{center}

\begin{caution}{This is the sound being edited, not a stored one}
These addresses reach the instrument's working copy --- the sound currently loaded for
editing. They are not a way to read or write one of the 32 stored sounds. For those, use
the \textsc{sound} bulk-dump category, which carries the whole sound memory at
\bytes{20 00 00}.
\end{caution}

\section{How a normal sound is laid out}
The \textsc{normal sound} area is \textbf{713 bytes}. It is one \textsc{general data}
block, then four identical \textsc{tone data} blocks, then four identical
\textsc{modeling data} blocks --- the instrument's four tones and the four driver and
resonator models that go with them.

\begin{center}
\begin{tabular}{lrrl}
\toprule
block & size & repeats & parameter numbers \\
\midrule
\textsc{general data}  & 217 & 1 & \bytes{000}--\bytes{0D8} \\
\textsc{tone data}     &  81 & 4 & \bytes{0D9}, \bytes{12A}, \bytes{17B}, \bytes{1CC} \\
\textsc{modeling data} &  43 & 4 & \bytes{21D}, \bytes{248}, \bytes{273}, \bytes{29E} \\
\midrule
total & 713 & & \bytes{000}--\bytes{2C8} \\
\bottomrule
\end{tabular}
\end{center}

A parameter number is an offset from the start of the area, so the address of parameter
\textit{p} of tone \textit{n} (counting tones from 1) is

\begin{center}
\bytes{040000} $+$ \bytes{0D9} $+$ 81 $\times$ (\textit{n} $-$ 1) $+$ \textit{p}
\end{center}

expressed as \textsc{adr} in the usual three seven-bit bytes, and the same with \bytes{21D}
and 43 for a modeling block. The tables below give one instance of each repeating block;
the other three are identical at the offsets above.

\begin{note}{Two independent sources agree on the size}
The block sizes and repeat counts here come from Technics' published parameter tables. The
figure of 713 bytes for the whole area comes from the program, decoded separately, and the
guide never states it. The transcribed blocks tile \bytes{000} to \bytes{2C8} with no gap
and no overlap and come to exactly that total, which is what makes the layout above
trustworthy rather than merely plausible.
\end{note}

\section{What is transcribed here, and what is not}
The tables give each parameter's number, its bit field within that byte, its name and its
range. They were read from scans of the published tables, and the numbering is checked by
the tiling described above. \textbf{The names, bit fields and ranges are not checked by
anything}; a misreading of one would survive.

Three enumerations the guide prints are not reproduced: the list of about fifty control
functions that a controller's \textsc{function select} byte chooses from, the meaning of
the filter \textsc{value} bytes for each filter mode, and the meaning of the twenty
\textsc{value} bytes under each effect type. All three are in the guide.

%% GENERATED by notes/sysex-probes/sysex_sound_layout.py --tex
%% Source: sound_layout.json. Do not edit.

\section{Normal Sound parameters}
\subsection*{GENERAL DATA}
{\small
\begin{longtable}{l>{\raggedright\arraybackslash}p{24mm}>{\raggedright\arraybackslash}p{58mm}>{\raggedright\arraybackslash}p{36mm}}
\toprule
no. & bits & parameter & values \\
\midrule\endfirsthead
\toprule no. & bits & parameter & values \\
\midrule\endhead
\bytes{000--00F} & 6--0 & NAME (16 characters) & \bytes{20-7F} --- space to \} \\
\bytes{010} & 7--6 & ATTRIBUTE FLAG: PARAMETER MODE & \bytes{00-00} --- 0: NORMAL MODE, fixed \\
\bytes{011} & 7--0 & TONE ON/OFF & two bits per tone, 4th 3rd 2nd 1st \\
\bytes{012} & 7--0 & KEY ON MODE & per tone: 0 KEY ON, 1 KEY OFF, 2 LEGATO, 3 NON LEGATO \\
\bytes{013} & 7--0 & KEY SCALING & \bytes{00-42} --- 0 OFF, 1-22 octave, 64-66 all key \\
\bytes{014--016} &  & CONTROLLERS: PITCH BEND & on/off and phase, function select, depth \\
\bytes{017--019} &  & MODULATION WHEEL 1 SELECT1 & same three bytes \\
\bytes{01A--01C} &  & MODULATION WHEEL 1 SELECT2 &  \\
\bytes{01D--01F} &  & MODULATION WHEEL 2 SELECT1 &  \\
\bytes{020--022} &  & MODULATION WHEEL 2 SELECT2 &  \\
\bytes{023} & 6--0 & REALTIME CREATOR BOX SELECT & bit per box, BOX1-BOX6 \\
\bytes{024} & 6--0 & REALTIME CONTROLLER BOX SELECT & bit per box, BOX1-BOX6 \\
\bytes{025--027} &  & CONTROL PEDAL SELECT1 &  \\
\bytes{028--02A} &  & CONTROL PEDAL SELECT2 &  \\
\bytes{02B--02D} &  & AFTER TOUCH SELECT1 &  \\
\bytes{02E--030} &  & AFTER TOUCH SELECT2 &  \\
\bytes{031--054} &  & BOX SELECT: CONTROLLER BOX1 X through BOX6 Y & twelve controllers, three bytes each, X then Y \\
\bytes{055} & 3--0 & PITCH OCTAVE SHIFT & \bytes{06-0A} --- 6 is -2 oct, 8 is 0, A is +2 oct \\
\bytes{056} &  & (RESERVE) &  \\
\bytes{057--066} &  & PITCH LFO BOX1-BOX4 & four bytes per box: depth, speed, key sync and touch depth, wave and delay \\
\bytes{067--076} &  & LEVEL LFO BOX1-BOX4 & same four bytes per box \\
\bytes{077--086} &  & FILTER LFO BOX1-BOX4 & same four bytes per box \\
\bytes{087} & 3--0, 7--4 & MIXER OUTPUT SELECT: MAIN, SUB & \bytes{00-05} --- 0 no output, 1 MAIN, 2-4 SUB1-3, 5 EFFECT2 \\
\bytes{088} & 6--0 & EFFECT SEND VOLUME: EFFECT1 & \bytes{00-7F} \\
\bytes{089} & 6--0 & EFFECT SEND VOLUME: REVERB & \bytes{00-7F} \\
\bytes{08A} & 3--0, 7--4 & EFFECT OUTPUT SELECT: EFFECT1, EFFECT2 & \bytes{01-04} --- 1 MAIN, 2-4 SUB1-3 \\
\bytes{08B} & 0 & DSP EFFECT COMMON ALGORITHM & \bytes{00-01} --- 0 SERIAL, 1 PARALLEL \\
\bytes{08C} &  & EFFECT1 TYPE & see the guide's DSP EFFECT pages \\
\bytes{08D--0A0} &  & EFFECT1 PARAMETER VALUE1-VALUE20 & meaning depends on the type \\
\bytes{0A1} &  & EFFECT1 ANOTHER EFF SEND &  \\
\bytes{0A2} &  & EFFECT1 DYNAMIC CONTROL &  \\
\bytes{0A3} &  & EFFECT2 TYPE &  \\
\bytes{0A4--0B7} &  & EFFECT2 PARAMETER VALUE1-VALUE20 &  \\
\bytes{0B8} &  & EFFECT2 ANOTHER EFF SEND &  \\
\bytes{0B9} &  & EFFECT2 DYNAMIC CONTROL &  \\
\bytes{0BA} &  & REVERB TYPE &  \\
\bytes{0BB--0CE} &  & REVERB PARAMETER VALUE1-VALUE20 &  \\
\bytes{0CF} &  & REVERB DYNAMIC CONTROL & reverb has no ANOTHER EFF SEND \\
\bytes{0D0} & 7, 6, 3--0 & DIGITAL EFFECT MODE: ON/OFF, PAN, EFFECT TYPE & \bytes{00-0B} --- 0-1 CELESTE, 2-3 CHORUS, 4-5 ENSEMBLE, 6-7 TREMOLO, 8-9 DELAY, 10-11 SOLO EFFECT \\
\bytes{0D1--0D8} &  & DIGITAL EFFECT VALUE1-VALUE8 &  \\
\bottomrule
\end{longtable}
}

\subsection*{TONE DATA}
{\small
\begin{longtable}{l>{\raggedright\arraybackslash}p{24mm}>{\raggedright\arraybackslash}p{58mm}>{\raggedright\arraybackslash}p{36mm}}
\toprule
no. & bits & parameter & values \\
\midrule\endfirsthead
\toprule no. & bits & parameter & values \\
\midrule\endhead
\bytes{0D9} & 5--0 & ATTRIBUTE DELAY & \bytes{00-32} --- 0 delay off \\
\bytes{0DA} & 7--0 & ATTRIBUTE PANNING & \bytes{00-80} --- 0-127 fixed pan, 128 random pan \\
\bytes{0DB} & 6--0 & TONE SELECT: TTN SELECT & \bytes{00-7F} \\
\bytes{0DC} & 7--6, 5--4, 3--0 & TONE KIND, TONE LOCATION, BANK SELECT & kind 0 NORMAL 1 DRUM 2 ATTACK; location 0 internal ROM 1 internal RAM 2 external ROM \\
\bytes{0DD} & 7--0 & PITCH KEY SHIFT & \bytes{E8-18} --- -24 to +24 \\
\bytes{0DE} & 7--0 & PITCH DETUNE & \bytes{80-7F} --- -128 to +127 \\
\bytes{0DF} & 7--6, 5, 4, 2--0 & BOX CONTROL LFO-FM1: BOX SELECT, ON/OFF, PHASE, SCALE SELECT & scale 0 1/1 through 6 1/64, 7 FIX \\
\bytes{0E0--0E9} &  & PITCH ENVELOPE SHAPE: total depth, start pitch, attack time, peak level, decay1 time, sustain1 level, decay2 time, sustain2 level, release time, stop pitch & one byte each, in that order \\
\bytes{0EA} & 7--0 & PITCH ENV TOUCH FOLLOW ADR TIME & \bytes{CE-32} --- -50 to +50 \\
\bytes{0EB} & 7--0 & PITCH ENV TOUCH FOLLOW DEPTH & \bytes{CE-32} \\
\bytes{0EC} & 6--0 & PITCH ENV KEY FOLLOW CENTER KEY & \bytes{00-7F} \\
\bytes{0ED} & 7--0 & PITCH ENV KEY FOLLOW ATTACK SLOPE & \bytes{CE-32} \\
\bytes{0EE} & 7--0 & PITCH ENV KEY FOLLOW DECAY SLOPE & \bytes{CE-32} \\
\bytes{0EF} & 7--0 & PITCH ENV KEY FOLLOW RELEASE SLOPE & \bytes{CE-32} \\
\bytes{0F0} & 6--0 & LEVEL VOLUME & \bytes{00-7F} --- one step is 0.375 dB \\
\bytes{0F1} & 7--0 & LEVEL TOUCH DEPTH & \bytes{CE-32} \\
\bytes{0F2} & 7--5 & LEVEL TOUCH CURVE & \bytes{00-06} --- 0 is -3 curve through 6 is +3 curve \\
\bytes{0F3} & 6--0 & LEVEL KEY FOLLOW CENTER KEY & \bytes{00-7F} \\
\bytes{0F4} & 6--0 & LEVEL KEY FOLLOW LOW KEY LIMIT & \bytes{00-7F} \\
\bytes{0F5} & 6--0 & LEVEL KEY FOLLOW HIGH KEY LIMIT & \bytes{00-7F} \\
\bytes{0F6} & 7--0 & LEVEL KEY FOLLOW SLOPE & \bytes{CE-32} \\
\bytes{0F7--0FA} & 6--0 & LAYER KEY RANGE: low key select, low fade, high key select, high fade & \bytes{00-7F} \\
\bytes{0FB--0FE} & 6--0 & LAYER VELOCITY RANGE: low vel select, low fade, high vel select, high fade & \bytes{00-7F} \\
\bytes{0FF} & 7--6, 5, 4 & BOX CONTROL LFO-AM: BOX SELECT, ON/OFF, PHASE &  \\
\bytes{100--106} &  & LEVEL ENVELOPE SHAPE: attack time, peak level, decay1 time, sustain1 level, decay2 time, sustain2 level, release time &  \\
\bytes{107} & 7--0 & LEVEL ENV TOUCH FOLLOW ATTACK & \bytes{CE-32} \\
\bytes{108} & 7--0 & LEVEL ENV TOUCH FOLLOW DECAY & \bytes{CE-32} \\
\bytes{109--10B} & 6--0 & LEVEL ENV KEY FOLLOW: center key, low key limit, high key limit & \bytes{00-7F} \\
\bytes{10C--10E} & 7--0 & LEVEL ENV KEY FOLLOW: attack slope, decay slope, release slope & \bytes{CE-32} \\
\bytes{10F} & 2--0, 7--5 & FILTER MODE, TOUCH CURVE & \bytes{00-06} --- 0 THROUGH, 1 LPF12+EQ, 2 HPF12+EQ, 3 LPF24, 4 HPF24, 5 BPF \\
\bytes{110} & 5--0 & FILTER TOUCH DEPTH & \bytes{00-32} \\
\bytes{111} & 7--6, 5, 4 & FILTER LFO-FCM: BOX SELECT, ON/OFF, PHASE &  \\
\bytes{112--114} & 6--0 & FILTER KEY FOLLOW: center key, low key limit, high key limit & \bytes{00-7F} \\
\bytes{115} & 7--0 & FILTER KEY FOLLOW SLOPE & \bytes{CE-32} \\
\bytes{116--11F} &  & FILTER ENVELOPE SHAPE: cutoff adjust, start point, attack time, peak level, decay1 time, sustain1 level, decay2 time, sustain2 level, release time, stop point &  \\
\bytes{120} & 7--0 & FILTER ENV TOUCH FOLLOW ADR TIME & \bytes{CE-32} \\
\bytes{121} & 7--0 & FILTER ENV TOUCH FOLLOW DEPTH & \bytes{CE-32} \\
\bytes{122} & 6--0 & FILTER ENV KEY FOLLOW CENTER KEY & \bytes{00-7F} \\
\bytes{123--125} & 7--0 & FILTER ENV KEY FOLLOW: attack slope, decay slope, release slope & \bytes{CE-32} \\
\bytes{126--129} &  & FILTER PARAMETER VALUE1-VALUE4 & meaning depends on the filter mode \\
\bottomrule
\end{longtable}
}

\subsection*{MODELING DATA}
{\small
\begin{longtable}{l>{\raggedright\arraybackslash}p{24mm}>{\raggedright\arraybackslash}p{58mm}>{\raggedright\arraybackslash}p{36mm}}
\toprule
no. & bits & parameter & values \\
\midrule\endfirsthead
\toprule no. & bits & parameter & values \\
\midrule\endhead
\bytes{21D--21F} & 6--0 & DRIVER THRESHOLD POINT: between A and B, B and C, C and D & \bytes{00-7F} --- provided that A < B < C < D \\
\bytes{220} & 6--0 & DRIVER A TTN SELECT & \bytes{00-7F} \\
\bytes{221} & 7--6, 5--4, 3--0 & DRIVER A KIND, LOCATION, BANK SELECT & kind 0 NORMAL 1 DRUM 2 ATTACK \\
\bytes{222--223} &  & DRIVER B: TTN select; kind, location, bank select &  \\
\bytes{224--225} &  & DRIVER C: TTN select; kind, location, bank select &  \\
\bytes{226--227} &  & DRIVER D: TTN select; kind, location, bank select &  \\
\bytes{228} & 7--6, 5--0 & RESONATOR GENERAL: INTERACTION CONNECT, RESONATOR TYPE & \bytes{00-3F} --- type 0 ORIGINAL, 1-63 preset \\
\bytes{229} &  & (RESERVE) &  \\
\bytes{22A} & 7--0 & POSITION PARAMETER POSITION & \bytes{00-FA} --- 0.2 per cent per step, 0 to 50 per cent \\
\bytes{22B} & 7, 6--0 & FORMANT FLAG, DEPTH & \bytes{00-7F} --- 0 FIXED, 1 MOVE \\
\bytes{22C} &  & (RESERVE) &  \\
\bytes{22D} & 7--0 & POSITION TOUCH & \bytes{CE-32} --- -50 to +50 \\
\bytes{22E} & 5--0 & POSITION MOVEMENT WIDTH & \bytes{00-32} \\
\bytes{22F} & 5--0, 7 & POSITION MOVEMENT SPEED, SAMPLE/HOLD & \bytes{00-32} \\
\bytes{230} & 6--0 & INTERACTION GAIN & \bytes{00-7F} \\
\bytes{231} &  & (RESERVE) &  \\
\bytes{232} & 7, 6--0 & MAIN RESONATOR MODE, FITTING & \bytes{00-7F} --- 0 positive feedback, 1 negative feedback \\
\bytes{233} & 7, 6--0 & MAIN RESONATOR SCALE, MUTING & \bytes{00-7F} \\
\bytes{234} & 7--0 & MAIN TOUCH FOLLOW TOUCH FITTING & \bytes{CE-32} \\
\bytes{235} & 7--0 & MAIN TOUCH FOLLOW TOUCH MUTING & \bytes{CE-32} \\
\bytes{236--238} & 6--0 & MAIN KEY FOLLOW: center key, low key limit, high key limit & \bytes{00-7F} \\
\bytes{239} & 7--0 & MAIN KEY FOLLOW SLOPE & \bytes{CE-32} \\
\bytes{23A} & 7--0 & MAIN PITCH KEY SHIFT & \bytes{C4-3C} --- -60 to +60 \\
\bytes{23B} & 7--0 & MAIN PITCH DETUNE & \bytes{80-7F} --- -128 to +127 \\
\bytes{23C} & 7, 6--0 & SUB RESONATOR MODE, FITTING & \bytes{00-7F} \\
\bytes{23D} & 7, 6--0 & SUB RESONATOR SCALE, MUTING & \bytes{00-7F} \\
\bytes{23E} & 7--0 & SUB GAIN & \bytes{9C-64} \\
\bytes{23F} & 7--0 & SUB TOUCH FOLLOW TOUCH FITTING & \bytes{CE-32} \\
\bytes{240} & 7--0 & SUB TOUCH FOLLOW TOUCH MUTING & \bytes{CE-32} \\
\bytes{241} & 7--0 & SUB TOUCH FOLLOW TOUCH SUB GAIN & \bytes{CE-32} \\
\bytes{242--244} & 6--0 & SUB KEY FOLLOW: center key, low key limit, high key limit & \bytes{00-7F} \\
\bytes{245} & 7--0 & SUB KEY FOLLOW SLOPE & \bytes{CE-32} \\
\bytes{246} & 7--0 & SUB PITCH KEY SHIFT & \bytes{C4-3C} \\
\bytes{247} & 7--0 & SUB PITCH DETUNE & \bytes{80-7F} \\
\bottomrule
\end{longtable}
}

\section{Drum Sound parameters}
\subsection*{KIT DATA}
{\small
\begin{longtable}{l>{\raggedright\arraybackslash}p{24mm}>{\raggedright\arraybackslash}p{58mm}>{\raggedright\arraybackslash}p{36mm}}
\toprule
no. & bits & parameter & values \\
\midrule\endfirsthead
\toprule no. & bits & parameter & values \\
\midrule\endhead
\bytes{000--00F} & 6--0 & NAME (16 characters) & \bytes{20-7F} --- space to \} \\
\bytes{010} & 7--6 & ATTRIBUTE FLAG: PARAMETER MODE & \bytes{00-00} --- 2: DRUM MODE, fixed \\
\bytes{011--013} &  & CONTROLLERS: PITCH BEND & on/off and phase, function select, depth \\
\bytes{014--016} &  & MODULATION WHEEL 1 SELECT1 & on/off and phase, function select, depth \\
\bytes{017--019} &  & MODULATION WHEEL 1 SELECT2 & on/off and phase, function select, depth \\
\bytes{01A--01C} &  & MODULATION WHEEL 2 SELECT1 & on/off and phase, function select, depth \\
\bytes{01D--01F} &  & MODULATION WHEEL 2 SELECT2 & on/off and phase, function select, depth \\
\bytes{020} & 6--0 & REALTIME CREATOR BOX SELECT & bit per box, BOX1-BOX6 \\
\bytes{021} & 6--0 & REALTIME CONTROLLER BOX SELECT & bit per box, BOX1-BOX6 \\
\bytes{022--024} &  & CONTROL PEDAL SELECT1 & on/off and phase, function select, depth \\
\bytes{025--027} &  & CONTROL PEDAL SELECT2 & on/off and phase, function select, depth \\
\bytes{028--02A} &  & AFTER TOUCH SELECT1 & on/off and phase, function select, depth \\
\bytes{02B--02D} &  & AFTER TOUCH SELECT2 & on/off and phase, function select, depth \\
\bytes{02E--051} &  & BOX SELECT: CONTROLLER BOX1 X through BOX6 Y & twelve controllers, three bytes each, X then Y \\
\bytes{052} & 3--0, 7--4 & MIXER EFFECT OUTPUT SELECT: EFFECT1, EFFECT2 & \bytes{01-04} --- 1 MAIN, 2-4 SUB1-3 \\
\bytes{053} & 0 & DSP EFFECT COMMON ALGORITHM & \bytes{00-01} --- 0 SERIAL, 1 PARALLEL \\
\bytes{054} &  & EFFECT1 TYPE & see the effect tables \\
\bytes{055--068} &  & EFFECT1 PARAMETER VALUE1-VALUE20 & meaning depends on the type \\
\bytes{069} &  & EFFECT1 ANOTHER EFF SEND &  \\
\bytes{06A} &  & EFFECT1 DYNAMIC CONTROL &  \\
\bytes{06B} &  & EFFECT2 TYPE &  \\
\bytes{06C--07F} &  & EFFECT2 PARAMETER VALUE1-VALUE20 &  \\
\bytes{080} &  & EFFECT2 ANOTHER EFF SEND &  \\
\bytes{081} &  & EFFECT2 DYNAMIC CONTROL &  \\
\bytes{082} &  & REVERB TYPE &  \\
\bytes{083--096} &  & REVERB PARAMETER VALUE1-VALUE20 &  \\
\bytes{097} &  & REVERB DYNAMIC CONTROL & reverb has no ANOTHER EFF SEND \\
\bottomrule
\end{longtable}
}

\subsection*{SOUND SELECT}
{\small
\begin{longtable}{l>{\raggedright\arraybackslash}p{24mm}>{\raggedright\arraybackslash}p{58mm}>{\raggedright\arraybackslash}p{36mm}}
\toprule
no. & bits & parameter & values \\
\midrule\endfirsthead
\toprule no. & bits & parameter & values \\
\midrule\endhead
\bytes{098--197} &  & SOUND SELECT, NOTE:00 to NOTE:7F & 128 notes, two bytes each \\
\bottomrule
\end{longtable}
}

\subsection*{NOTE GENERAL DATA}
{\small
\begin{longtable}{l>{\raggedright\arraybackslash}p{24mm}>{\raggedright\arraybackslash}p{58mm}>{\raggedright\arraybackslash}p{36mm}}
\toprule
no. & bits & parameter & values \\
\midrule\endfirsthead
\toprule no. & bits & parameter & values \\
\midrule\endhead
\bytes{198--1A4} & 6--0 & NAME (13 characters) & \bytes{20-7F} \\
\bytes{1A5} & 5, 3--0 & ATTRIBUTE FLAG: KEY OFF MODE, TONE ON/OFF & \bytes{00-01} --- two bits per tone, 2nd then 1st \\
\bytes{1A6} & 7, 6--0 & SAME GROUP DUMP: DUMP FLAG, GROUP NUMBER & \bytes{00-7F} --- flag 1 dumps the same group's sound \\
\bytes{1A7} & 3--0, 7--4 & MIXER OUTPUT SELECT: MAIN, SUB & \bytes{00-05} --- 0 no output, 1 MAIN, 2-4 SUB1-3, 5 EFFECT2 \\
\bytes{1A8} & 6--0 & EFFECT SEND VOLUME: EFFECT1 & \bytes{00-7F} \\
\bytes{1A9} & 6--0 & EFFECT SEND VOLUME: REVERB & \bytes{00-7F} \\
\bottomrule
\end{longtable}
}

\subsection*{NOTE TONE DATA}
{\small
\begin{longtable}{l>{\raggedright\arraybackslash}p{24mm}>{\raggedright\arraybackslash}p{58mm}>{\raggedright\arraybackslash}p{36mm}}
\toprule
no. & bits & parameter & values \\
\midrule\endfirsthead
\toprule no. & bits & parameter & values \\
\midrule\endhead
\bytes{1AA} & 5--0 & ATTRIBUTE DELAY & \bytes{00-32} --- 0 delay off \\
\bytes{1AB} & 7--0 & ATTRIBUTE PANNING & \bytes{00-80} --- 0-127 fixed pan, 128 random pan \\
\bytes{1AC} & 6--0 & TONE SELECT: TTN SELECT & \bytes{00-7F} \\
\bytes{1AD} & 7--6, 5--4, 3--0 & TONE KIND, TONE LOCATION, BANK SELECT & kind 0 NORMAL 1 DRUM 2 ATTACK \\
\bytes{1AE} & 7--0 & PITCH KEY SHIFT & \bytes{E8-18} --- -24 to +24 \\
\bytes{1AF} & 7--0 & PITCH DETUNE & \bytes{80-7F} --- -128 to +127 \\
\bytes{1B0} & 6--0 & LEVEL VOLUME & \bytes{00-7F} --- one step is 0.375 dB \\
\bytes{1B1} & 7--0 & LEVEL TOUCH DEPTH & \bytes{CE-32} --- -50 to +50 \\
\bytes{1B2} & 7--5 & LEVEL TOUCH CURVE & \bytes{00-06} --- 0 is -3 curve through 6 is +3 curve \\
\bytes{1B3--1B8} & 6--0 & LEVEL ENVELOPE SHAPE: attack time, decay1 time, sustain1 level, decay2 time, sustain2 level, release time & \bytes{00-64} --- one byte each, in that order \\
\bytes{1B9} & 7--0 & LEVEL ENV TOUCH FOLLOW ATTACK & \bytes{CE-32} \\
\bytes{1BA} & 7--0 & LEVEL ENV TOUCH FOLLOW DECAY & \bytes{CE-32} \\
\bytes{1BB} & 2--0, 7--5 & FILTER MODE, TOUCH CURVE & \bytes{00-06} --- 0 THROUGH, 1 LPF12+EQ, 2 HPF12+EQ, 3 LPF24, 4 HPF24, 5 BPF \\
\bytes{1BC} & 5--0 & FILTER TOUCH DEPTH & \bytes{00-32} \\
\bytes{1BD--1C0} &  & FILTER PARAMETER VALUE1-VALUE4 & meaning depends on the filter mode \\
\bottomrule
\end{longtable}
}

\subsection*{NOTE MODELING DATA}
{\small
\begin{longtable}{l>{\raggedright\arraybackslash}p{24mm}>{\raggedright\arraybackslash}p{58mm}>{\raggedright\arraybackslash}p{36mm}}
\toprule
no. & bits & parameter & values \\
\midrule\endfirsthead
\toprule no. & bits & parameter & values \\
\midrule\endhead
\bytes{1D8--1DA} & 6--0 & DRIVER THRESHOLD POINT: between A and B, B and C, C and D & \bytes{00-7F} --- provided that A < B < C < D \\
\bytes{1DB} & 6--0 & DRIVER A TTN SELECT & \bytes{00-7F} \\
\bytes{1DC} & 7--6, 5--4, 3--0 & DRIVER A KIND, LOCATION, BANK SELECT & kind 0 NORMAL 1 DRUM 2 ATTACK \\
\bytes{1DD--1DE} &  & DRIVER B: TTN select; kind, location, bank select &  \\
\bytes{1DF--1E0} &  & DRIVER C: TTN select; kind, location, bank select &  \\
\bytes{1E1--1E2} &  & DRIVER D: TTN select; kind, location, bank select &  \\
\bytes{1E3} & 7--6, 5--0 & RESONATOR GENERAL: INTERACTION CONNECT, RESONATOR TYPE & \bytes{00-3F} --- type 0 ORIGINAL, 1-63 preset \\
\bytes{1E4} &  & (RESERVE) &  \\
\bytes{1E5} & 7--0 & POSITION PARAMETER POSITION & \bytes{00-FA} --- 0.2 per cent per step, 0 to 50 per cent \\
\bytes{1E6} & 7, 6--0 & FORMANT FLAG, DEPTH & \bytes{00-7F} --- 0 FIXED, 1 MOVE \\
\bytes{1E7} &  & (RESERVE) &  \\
\bytes{1E8} & 7--0 & POSITION TOUCH & \bytes{CE-32} --- -50 to +50 \\
\bytes{1E9} & 5--0 & POSITION MOVEMENT WIDTH & \bytes{00-32} \\
\bytes{1EA} & 5--0, 7 & POSITION MOVEMENT SPEED, SAMPLE/HOLD & \bytes{00-32} \\
\bytes{1EB} & 6--0 & INTERACTION GAIN & \bytes{00-7F} \\
\bytes{1EC} &  & (RESERVE) &  \\
\bytes{1ED} & 7, 6--0 & MAIN RESONATOR MODE, FITTING & \bytes{00-7F} --- 0 positive feedback, 1 negative feedback \\
\bytes{1EE} & 7, 6--0 & MAIN RESONATOR SCALE, MUTING & \bytes{00-7F} \\
\bytes{1EF} & 7--0 & MAIN TOUCH FOLLOW TOUCH FITTING & \bytes{CE-32} \\
\bytes{1F0} & 7--0 & MAIN TOUCH FOLLOW TOUCH MUTING & \bytes{CE-32} \\
\bytes{1F1--1F3} & 6--0 & MAIN KEY FOLLOW: center key, low key limit, high key limit & \bytes{00-7F} \\
\bytes{1F4} & 7--0 & MAIN KEY FOLLOW SLOPE & \bytes{CE-32} \\
\bytes{1F5} & 7--0 & MAIN PITCH KEY SHIFT & \bytes{C4-3C} --- -60 to +60 \\
\bytes{1F6} & 7--0 & MAIN PITCH DETUNE & \bytes{80-7F} --- -128 to +127 \\
\bytes{1F7} & 7, 6--0 & SUB RESONATOR MODE, FITTING & \bytes{00-7F} \\
\bytes{1F8} & 7, 6--0 & SUB RESONATOR SCALE, MUTING & \bytes{00-7F} \\
\bytes{1F9} & 7--0 & SUB GAIN & \bytes{9C-64} \\
\bytes{1FA} & 7--0 & SUB TOUCH FOLLOW TOUCH FITTING & \bytes{CE-32} \\
\bytes{1FB} & 7--0 & SUB TOUCH FOLLOW TOUCH MUTING & \bytes{CE-32} \\
\bytes{1FC} & 7--0 & SUB TOUCH FOLLOW TOUCH SUB GAIN & \bytes{CE-32} \\
\bytes{1FD--1FF} & 6--0 & SUB KEY FOLLOW: center key, low key limit, high key limit & \bytes{00-7F} \\
\bytes{200} & 7--0 & SUB KEY FOLLOW SLOPE & \bytes{CE-32} \\
\bytes{201} & 7--0 & SUB PITCH KEY SHIFT & \bytes{C4-3C} \\
\bytes{202} & 7--0 & SUB PITCH DETUNE & \bytes{80-7F} \\
\bottomrule
\end{longtable}
}



\section{Drum sound}
\textsc{drum sound} is reached at \bytes{18 00 00}, and messages addressed there are
accepted. Its contents are not enumerated in this edition.

\chapter{Effects}\label{ch:effects}

A sound carries three effect blocks --- \textsc{effect1}, \textsc{effect2} and
\textsc{reverb} --- and each is a \textsc{type} byte followed by twenty \textsc{value}
bytes, an \textsc{another eff send} and a \textsc{dynamic control}. The \textsc{type} byte
selects one of 56 effects, and what the twenty \textsc{value} bytes mean depends entirely
on which.

\section{What the type byte selects}
The numbers are not contiguous: they are grouped, with modulation effects below
\bytes{10}, reverbs at \bytes{10}--\bytes{1B}, distortion and dynamics at
\bytes{20}--\bytes{27}, a second modulation group at \bytes{30}--\bytes{38}, two-effect
combinations at \bytes{40}--\bytes{4B} and three-effect combinations at
\bytes{60}--\bytes{63}.

Twelve of the reverbs are marked \textsc{(rev only)} in the published guide: they are
intended for the \textsc{reverb} block rather than \textsc{effect1} or \textsc{effect2}.

\begin{note}{Two sources agree on which effects exist}
The table below was transcribed from the guide. Independently, this project
disassembled the effects DSP and produced a manifest of every distinct effect program
with the number that selects it. The two sets of numbers are equal --- 56 effects,
no more on either side --- which is what makes the numbering here trustworthy.

They disagree only in spelling, because the ROM carries the short names the instrument
displays and the guide prints long ones. The table gives the guide's.
\end{note}

\section{What the value bytes mean}
The \textsc{value} list below is the order the bytes appear in, so the first name is
\textsc{value1}, the second \textsc{value2}, and so on. An effect that names fewer than
twenty leaves the rest unused.

The combination effects concatenate their components: the first block's parameters, then
a dry/wet for the second block, then the second block's. A three-effect combination
continues the same way.

\begin{caution}{The value lists are transcribed, and not checked by anything}
The effect \emph{numbers} are confirmed twice over. The \emph{names} in the value lists
are not: they were read from scans, nothing independent states them, and a misreading
would survive. Treat them as a guide to what the bytes are for rather than as a
specification.

Two further things the guide gives and this edition does not reproduce: the conversion
from a byte value to a physical unit --- seconds, hertz, decibels --- which it prints as
eighteen tables, and the block diagram of each effect.
\end{caution}

%% GENERATED by notes/sysex-probes/sysex_effects.py --tex
%% Source: effects.json, effect_value_numbers.json and
%% dsp/programs.tsv. Do not edit.

{\small
\begin{longtable}{ll>{\raggedright\arraybackslash}p{40mm}>{\raggedright\arraybackslash}p{78mm}}
\toprule
no. & n & effect & parameters, with the \textsc{value} byte each occupies \\
\midrule\endfirsthead
\toprule no. & n & effect & parameters, with the \textsc{value} byte each occupies \\
\midrule\endhead
\bytes{00} & 0 & NO OPERATION & --- \\
\bytes{01} & 5 & CHORUS & WET~\textbf{1}, DEPTH~\textbf{2}, LFO SPEED~\textbf{3}, LFO WAVEFORM~\textbf{4}, VOLUME~\textbf{5} \\
\bytes{02} & 7 & MODULATED CHORUS & WET~\textbf{1}, DEPTH~\textbf{2}, SLOW LFO SPEED~\textbf{3}, FAST LFO SPEED~\textbf{4}, FAST LFO BALANCE~\textbf{5}, LFO WAVEFORM~\textbf{6}, VOLUME~\textbf{7} \\
\bytes{03} & 9 & ENHANCER & WET~\textbf{1}, MANUAL~\textbf{2}, LOW MIX~\textbf{3}, HIGH MIX~\textbf{4}, DELAY TIME L~\textbf{5--6}, DELAY TIME R~\textbf{7--8}, VOLUME~\textbf{9} \\
\bytes{04} & 8 & FLANGER & WET~\textbf{1}, DEPTH~\textbf{2}, LFO SPEED~\textbf{3}, RESONANCE~\textbf{4}, MANUAL~\textbf{5}, PHASE~\textbf{6}, LFO WAVEFORM~\textbf{7}, VOLUME~\textbf{8} \\
\bytes{05} & 8 & PHASER & WET~\textbf{1}, DEPTH~\textbf{2}, LFO SPEED~\textbf{3}, RESONANCE~\textbf{4}, MANUAL~\textbf{5}, PHASE~\textbf{6}, LFO WAVEFORM~\textbf{7}, VOLUME~\textbf{8} \\
\bytes{06} & 5 & ENSEMBLE & WET~\textbf{1}, DEPTH~\textbf{2}, LFO SPEED~\textbf{3}, LFO WAVEFORM~\textbf{4}, VOLUME~\textbf{5} \\
\bytes{08} & 6 & GATED REVERB & WET~\textbf{1}, GATE TIME~\textbf{2}, HIGH DUMP GAIN~\textbf{3}, THRESHOLD~\textbf{4}, MASK TIME~\textbf{5}, VOLUME~\textbf{6} \\
\bytes{09} & 9 & SINGLE DELAY & WET~\textbf{1}, DELAY L~\textbf{2--3}, DELAY R~\textbf{4--5}, FEEDBACK L~\textbf{6}, FEEDBACK R~\textbf{7}, HIGH DUMP GAIN~\textbf{8}, VOLUME~\textbf{9} \\
\bytes{0A} & 16 & MULTI TAP DELAY & WET~\textbf{1}, DELAY 1~\textbf{2--3}, DELAY 2~\textbf{4--5}, DELAY 3~\textbf{6--7}, DELAY 4~\textbf{8--9}, PAN 1~\textbf{10}, PAN 2~\textbf{11}, PAN 3~\textbf{12}, PAN 4~\textbf{13}, FEED BACK~\textbf{14}, HIGH DUMP GAIN~\textbf{15}, VOLUME~\textbf{16} \\
\bytes{0B} & 10 & MANUAL DELAY & WET~\textbf{1}, MODULATION DEPTH~\textbf{2}, DELAY L~\textbf{3--4}, DELAY R~\textbf{5--6}, FEEDBACK L~\textbf{7}, FEEDBACK R~\textbf{8}, HIGH DUMP GAIN~\textbf{9}, VOLUME~\textbf{10} \\
\bytes{10} & 5 & ROOM REVERB 1 (REV only) & REVERB TIME~\textbf{1}, PRE DELAY~\textbf{2}, HIGH DUMP GAIN~\textbf{3}, EARLY REFL LEVEL~\textbf{4}, VOLUME~\textbf{5} \\
\bytes{11} & 5 & ROOM REVERB 2 (REV only) & REVERB TIME~\textbf{1}, PRE DELAY~\textbf{2}, HIGH DUMP GAIN~\textbf{3}, EARLY REFL LEVEL~\textbf{4}, VOLUME~\textbf{5} \\
\bytes{12} & 5 & PLATE REVERB 1 (REV only) & REVERB TIME~\textbf{1}, PRE DELAY~\textbf{2}, HIGH DUMP GAIN~\textbf{3}, EARLY REFL LEVEL~\textbf{4}, VOLUME~\textbf{5} \\
\bytes{13} & 5 & PLATE REVERB 2 (REV only) & REVERB TIME~\textbf{1}, PRE DELAY~\textbf{2}, HIGH DUMP GAIN~\textbf{3}, EARLY REFL LEVEL~\textbf{4}, VOLUME~\textbf{5} \\
\bytes{14} & 5 & CONCERT REVERB 1 (REV only) & REVERB TIME~\textbf{1}, PRE DELAY~\textbf{2}, HIGH DUMP GAIN~\textbf{3}, EARLY REFL LEVEL~\textbf{4}, VOLUME~\textbf{5} \\
\bytes{15} & 5 & CONCERT REVERB 2 (REV only) & REVERB TIME~\textbf{1}, PRE DELAY~\textbf{2}, HIGH DUMP GAIN~\textbf{3}, EARLY REFL LEVEL~\textbf{4}, VOLUME~\textbf{5} \\
\bytes{16} & 5 & DARK REVERB 1 (REV only) & REVERB TIME~\textbf{1}, PRE DELAY~\textbf{2}, HIGH DUMP GAIN~\textbf{3}, EARLY REFL LEVEL~\textbf{4}, VOLUME~\textbf{5} \\
\bytes{17} & 5 & DARK REVERB 2 (REV only) & REVERB TIME~\textbf{1}, PRE DELAY~\textbf{2}, HIGH DUMP GAIN~\textbf{3}, EARLY REFL LEVEL~\textbf{4}, VOLUME~\textbf{5} \\
\bytes{18} & 5 & BRIGHT REVERB 1 (REV only) & REVERB TIME~\textbf{1}, PRE DELAY~\textbf{2}, HIGH DUMP GAIN~\textbf{3}, EARLY REFL LEVEL~\textbf{4}, VOLUME~\textbf{5} \\
\bytes{19} & 5 & BRIGHT REVERB 2 (REV only) & REVERB TIME~\textbf{1}, PRE DELAY~\textbf{2}, HIGH DUMP GAIN~\textbf{3}, EARLY REFL LEVEL~\textbf{4}, VOLUME~\textbf{5} \\
\bytes{1A} & 5 & WAVE REVERB 1 (REV only) & REVERB TIME~\textbf{1}, PRE DELAY~\textbf{2}, HIGH DUMP GAIN~\textbf{3}, EARLY REFL LEVEL~\textbf{4}, VOLUME~\textbf{5} \\
\bytes{1B} & 5 & WAVE REVERB 2 (REV only) & REVERB TIME~\textbf{1}, PRE DELAY~\textbf{2}, HIGH DUMP GAIN~\textbf{3}, EARLY REFL LEVEL~\textbf{4}, VOLUME~\textbf{5} \\
\bytes{20} & 4 & DISTORTION & WET~\textbf{1}, DRIVE~\textbf{2}, ADJUST~\textbf{3}, VOLUME~\textbf{4} \\
\bytes{21} & 4 & OVERDRIVE & WET~\textbf{1}, DRIVE~\textbf{2}, ADJUST~\textbf{3}, VOLUME~\textbf{4} \\
\bytes{22} & 4 & FUZZ & WET~\textbf{1}, DRIVE~\textbf{2}, ADJUST~\textbf{3}, VOLUME~\textbf{4} \\
\bytes{23} & 7 & EXCITER & WET~\textbf{1}, DRIVE~\textbf{2}, ADJUST~\textbf{3}, EMPHASIS Fc~\textbf{4--5}, EMPHASIS GAIN~\textbf{6}, VOLUME~\textbf{7} \\
\bytes{24} & 6 & COMPRESSOR & WET~\textbf{1}, THRESHOLD~\textbf{2}, RATIO~\textbf{3}, ATTACK SENSITIVITY~\textbf{4}, RELEASE SENSITIVITY~\textbf{5}, VOLUME~\textbf{6} \\
\bytes{25} & 5 & SLOW ATTACKER & WET~\textbf{1}, THRESHOLD~\textbf{2}, ATTACK RATE~\textbf{3}, RELEASE RATE~\textbf{4}, VOLUME~\textbf{5} \\
\bytes{26} & 1 & NOISE GENERATOR & VOLUME~\textbf{1} \\
\bytes{27} & 13 & PARAMETRIC EQ & BAND EMPHASIS 1 Fc~\textbf{1--2}, BAND EMPHASIS 1 Q~\textbf{1--2}, BAND EMPHASIS 1 G~\textbf{1--2}, BAND EMPHASIS 2 Fc~\textbf{3--4}, BAND EMPHASIS 2 Q~\textbf{3--4}, BAND EMPHASIS 2 G~\textbf{3--4}, BAND EMPHASIS 3 Fc~\textbf{5--6}, BAND EMPHASIS 3 Q~\textbf{5--6}, BAND EMPHASIS 3 G~\textbf{5--6}, BAND EMPHASIS 4 Fc~\textbf{7--8}, BAND EMPHASIS 4 Q~\textbf{7--8}, BAND EMPHASIS 4 G~\textbf{7--8}, BAND EMPHASIS 5 Fc~\textbf{9--10}, BAND EMPHASIS 5 Q~\textbf{9--10}, BAND EMPHASIS 5 G~\textbf{9--10}, BAND EMPHASIS 6 Fc~\textbf{11--12}, BAND EMPHASIS 6 Q~\textbf{11--12}, BAND EMPHASIS 6 G~\textbf{11--12}, VOLUME~\textbf{13} \\
\bytes{30} & 6 & AUTO PAN & WET~\textbf{1}, DEPTH~\textbf{2}, LFO SPEED~\textbf{3}, PHASE~\textbf{4}, LFO WAVEFORM~\textbf{5}, VOLUME~\textbf{6} \\
\bytes{31} & 6 & PITCH SHIFTER & WET~\textbf{1}, PITCH L~\textbf{2}, PITCH R~\textbf{3}, PRE DELAY~\textbf{4}, FEEDBACK~\textbf{5}, VOLUME~\textbf{6} \\
\bytes{32} & 6 & VIBRATO & WET~\textbf{1}, DEPTH~\textbf{2}, LFO SPEED~\textbf{3}, PHASE~\textbf{4}, LFO WAVEFORM~\textbf{5}, VOLUME~\textbf{6} \\
\bytes{33} & 6 & PEDAL WAH & WET~\textbf{1}, RESONANCE~\textbf{2}, MANUAL~\textbf{3}, SWEEP RANGE~\textbf{4}, WAH CENTER Fc~\textbf{5}, VOLUME~\textbf{6} \\
\bytes{34} & 5 & AUTO WAH & WET~\textbf{1}, RESONANCE~\textbf{2}, MANUAL~\textbf{3}, SWEEP RANGE~\textbf{4}, VOLUME~\textbf{5} \\
\bytes{35} & 15 & ROTARY SPEAKER & WET~\textbf{1}, DRIVE~\textbf{2}, VOLUME ADJUST~\textbf{3}, TREBLE DEPTH~\textbf{4}, TREBLE FAST~\textbf{5}, TREBLE SLOW~\textbf{6}, TREBLE WIND UP~\textbf{7}, TREBLE WIND DOWN~\textbf{8}, BASS DEPTH~\textbf{9}, BASS FAST~\textbf{10}, BASS SLOW~\textbf{11}, BASS WIND UP~\textbf{12}, BASS WIND DOWN~\textbf{13}, VOLUME~\textbf{14}, SLOW/FAST~\textbf{15} \\
\bytes{36} & 5 & RING MODULATOR & WET~\textbf{1}, OSC SPEED~\textbf{2}, PHASE~\textbf{3}, OSC WAVEFORM~\textbf{4}, VOLUME~\textbf{5} \\
\bytes{37} & 8 & HAAS EFFECT & WET~\textbf{1}, DELAY TIME L~\textbf{2--3}, DELAY TIME R~\textbf{4--5}, BALANCE L~\textbf{6}, BALANCE R~\textbf{7}, VOLUME~\textbf{8} \\
\bytes{38} & 8 & MIX UP & WET~\textbf{1}, DEPTH~\textbf{2}, SLOW LFO SPEED~\textbf{3}, FAST LFO SPEED L~\textbf{4}, FAST LFO SPEED R~\textbf{5}, PHASE~\textbf{6}, LFO WAVEFORM~\textbf{7}, VOLUME~\textbf{8} \\
\bytes{40} & 12 & SINGLE DELAY + CHORUS & DELAY WET~\textbf{1}, DELAY L~\textbf{2--3}, DELAY R~\textbf{4--5}, FEEDBACK L~\textbf{6}, FEEDBACK R~\textbf{7}, CHORUS DRY/WET~\textbf{8}, DEPTH~\textbf{9}, LFO SPEED~\textbf{10}, LFO WAVEFORM~\textbf{11}, VOLUME~\textbf{12} \\
\bytes{41} & 11 & SINGLE DELAY + SINGLE DELAY & DELAY 1 WET~\textbf{1}, DELAY L~\textbf{2}, DELAY R~\textbf{3}, FEEDBACK L~\textbf{4}, FEEDBACK R~\textbf{5}, DELAY 2 DRY/WET~\textbf{6}, DELAY L~\textbf{7}, DELAY R~\textbf{8}, FEEDBACK L~\textbf{9}, FEEDBACK R~\textbf{10}, VOLUME~\textbf{11} \\
\bytes{42} & 15 & SINGLE DELAY + FLANGER & DELAY WET~\textbf{1}, DELAY L~\textbf{2--3}, DELAY R~\textbf{4--5}, FEEDBACK L~\textbf{6}, FEEDBACK R~\textbf{7}, FLANGER DRY/WET~\textbf{8}, DEPTH~\textbf{9}, LFO SPEED~\textbf{10}, RESONANCE~\textbf{11}, MANUAL~\textbf{12}, PHASE~\textbf{13}, LFO WAVEFORM~\textbf{14}, VOLUME~\textbf{15} \\
\bytes{43} & 13 & SINGLE DELAY + VIBRATO & DELAY WET~\textbf{1}, DELAY L~\textbf{2--3}, DELAY R~\textbf{4--5}, FEEDBACK L~\textbf{6}, FEEDBACK R~\textbf{7}, VIBRATO DRY/WET~\textbf{8}, DEPTH~\textbf{9}, LFO SPEED~\textbf{10}, PHASE~\textbf{11}, LFO WAVEFORM~\textbf{12}, VOLUME~\textbf{13} \\
\bytes{44} & 15 & SINGLE DELAY + PHASER & DELAY WET~\textbf{1}, DELAY L~\textbf{2--3}, DELAY R~\textbf{4--5}, FEEDBACK L~\textbf{6}, FEEDBACK R~\textbf{7}, PHASER DRY/WET~\textbf{8}, DEPTH~\textbf{9}, LFO SPEED~\textbf{10}, RESONANCE~\textbf{11}, MANUAL~\textbf{12}, PHASE~\textbf{13}, LFO WAVEFORM~\textbf{14}, VOLUME~\textbf{15} \\
\bytes{45} & 13 & PEDAL WAH + DELAY & WAH WET~\textbf{1}, RESONANCE~\textbf{2}, MANUAL~\textbf{3}, SWEEP RANGE~\textbf{4}, WAH CENTER Fc~\textbf{5}, DELAY DRY/WET~\textbf{6}, DELAY L~\textbf{7--8}, DELAY R~\textbf{9--10}, FEEDBACK L~\textbf{11}, FEEDBACK R~\textbf{12}, VOLUME~\textbf{13} \\
\bytes{46} & 12 & AUTO WAH + SINGLE DELAY & WAH WET~\textbf{1}, RESONANCE~\textbf{2}, MANUAL~\textbf{3}, SWEEP RANGE~\textbf{4}, DELAY DRY/WET~\textbf{5}, DELAY L~\textbf{6--7}, DELAY R~\textbf{8--9}, FEEDBACK L~\textbf{10}, FEEDBACK R~\textbf{11}, VOLUME~\textbf{12} \\
\bytes{47} & 9 & PEQ + CHORUS & BAND EMPHASIS 1 Fc~\textbf{1--2}, BAND EMPHASIS 1 Q~\textbf{1--2}, BAND EMPHASIS 1 G~\textbf{1--2}, BAND EMPHASIS 2 Fc~\textbf{3--4}, BAND EMPHASIS 2 Q~\textbf{3--4}, BAND EMPHASIS 2 G~\textbf{3--4}, CHORUS DRY/WET~\textbf{5}, DEPTH~\textbf{6}, LFO SPEED~\textbf{7}, LFO WAVEFORM~\textbf{8}, VOLUME~\textbf{9} \\
\bytes{48} & 12 & PEQ + SINGLE DELAY & BAND EMPHASIS 1 Fc~\textbf{1--2}, BAND EMPHASIS 1 Q~\textbf{1--2}, BAND EMPHASIS 1 G~\textbf{1--2}, BAND EMPHASIS 2 Fc~\textbf{3--4}, BAND EMPHASIS 2 Q~\textbf{3--4}, BAND EMPHASIS 2 G~\textbf{3--4}, DELAY DRY/WET~\textbf{5}, DELAY L~\textbf{6--7}, DELAY R~\textbf{8--9}, FEEDBACK L~\textbf{10}, FEEDBACK R~\textbf{11}, VOLUME~\textbf{12} \\
\bytes{49} & 12 & PEQ + FLANGER & BAND EMPHASIS 1 Fc~\textbf{1--2}, BAND EMPHASIS 1 Q~\textbf{1--2}, BAND EMPHASIS 1 G~\textbf{1--2}, BAND EMPHASIS 2 Fc~\textbf{3--4}, BAND EMPHASIS 2 Q~\textbf{3--4}, BAND EMPHASIS 2 G~\textbf{3--4}, FLANGER DRY/WET~\textbf{5}, DEPTH~\textbf{6}, LFO SPEED~\textbf{7}, RESONANCE~\textbf{8}, MANUAL~\textbf{9}, PHASE~\textbf{10}, LFO WAVEFORM~\textbf{11}, VOLUME~\textbf{12} \\
\bytes{4A} & 10 & PEQ + VIBRATO & BAND EMPHASIS 1 Fc~\textbf{1--2}, BAND EMPHASIS 1 Q~\textbf{1--2}, BAND EMPHASIS 1 G~\textbf{1--2}, BAND EMPHASIS 2 Fc~\textbf{3--4}, BAND EMPHASIS 2 Q~\textbf{3--4}, BAND EMPHASIS 2 G~\textbf{3--4}, VIBRATO DRY/WET~\textbf{5}, DEPTH~\textbf{6}, LFO SPEED~\textbf{7}, PHASE~\textbf{8}, LFO WAVEFORM~\textbf{9}, VOLUME~\textbf{10} \\
\bytes{4B} & 9 & PEQ + COMPRESSOR & BAND EMPHASIS 1 Fc~\textbf{1--2}, BAND EMPHASIS 1 Q~\textbf{1--2}, BAND EMPHASIS 1 G~\textbf{1--2}, BAND EMPHASIS 2 Fc~\textbf{3--4}, BAND EMPHASIS 2 Q~\textbf{3--4}, BAND EMPHASIS 2 G~\textbf{3--4}, THRESHOLD~\textbf{5}, RATIO~\textbf{6}, ATTACK SENSITIVITY~\textbf{7}, RELEASE SENSITIVITY~\textbf{8}, VOLUME~\textbf{9} \\
\bytes{60} & 11 & PEQ + COMPRESSOR + DISTORTION & BAND EMPHASIS 1 Fc~\textbf{1--2}, BAND EMPHASIS 1 Q~\textbf{1--2}, BAND EMPHASIS 1 G~\textbf{1--2}, BAND EMPHASIS 2 Fc~\textbf{3--4}, BAND EMPHASIS 2 Q~\textbf{3--4}, BAND EMPHASIS 2 G~\textbf{3--4}, THRESHOLD~\textbf{5}, RATIO~\textbf{6}, ATTACK SENSITIVITY~\textbf{7}, RELEASE SENSITIVITY~\textbf{8}, DRIVE~\textbf{9}, ADJUST~\textbf{10}, VOLUME~\textbf{11} \\
\bytes{61} & 9 & PEQ + COMPRESSOR + OVERDRIVE & BAND EMPHASIS Fc~\textbf{1--2}, BAND EMPHASIS Q~\textbf{1--2}, BAND EMPHASIS G~\textbf{1--2}, THRESHOLD~\textbf{3}, RATIO~\textbf{4}, ATTACK SENSITIVITY~\textbf{5}, RELEASE SENSITIVITY~\textbf{6}, DRIVE~\textbf{7}, ADJUST~\textbf{8}, VOLUME~\textbf{9} \\
\bytes{62} & 14 & PEQ + DISTORTION + DELAY & BAND EMPHASIS 1 Fc~\textbf{1--2}, BAND EMPHASIS 1 Q~\textbf{1--2}, BAND EMPHASIS 1 G~\textbf{1--2}, BAND EMPHASIS 2 Fc~\textbf{3--4}, BAND EMPHASIS 2 Q~\textbf{3--4}, BAND EMPHASIS 2 G~\textbf{3--4}, DRIVE~\textbf{5}, ADJUST~\textbf{6}, DELAY DRY/WET~\textbf{7}, DELAY L~\textbf{8--9}, DELAY R~\textbf{10--11}, FEEDBACK L~\textbf{12}, FEEDBACK R~\textbf{13}, VOLUME~\textbf{14} \\
\bytes{63} & 14 & PEQ + OVERDRIVE + DELAY & BAND EMPHASIS 1 Fc~\textbf{1--2}, BAND EMPHASIS 1 Q~\textbf{1--2}, BAND EMPHASIS 1 G~\textbf{1--2}, BAND EMPHASIS 2 Fc~\textbf{3--4}, BAND EMPHASIS 2 Q~\textbf{3--4}, BAND EMPHASIS 2 G~\textbf{3--4}, DRIVE~\textbf{5}, ADJUST~\textbf{6}, DELAY DRY/WET~\textbf{7}, DELAY L~\textbf{8--9}, DELAY R~\textbf{10--11}, FEEDBACK L~\textbf{12}, FEEDBACK R~\textbf{13}, VOLUME~\textbf{14} \\
\bottomrule
\end{longtable}
}



\chapter{Enumerations}\label{ch:enums}

Three lists in chapter~\ref{ch:sound} are referred to rather than spelled out where they
are used, because each is shared by many parameters. They are given here.

\section{Control functions}\label{sec:controlfn}
Every controller in a sound --- the pitch bend lever, the two modulation wheels, the
control pedal, after touch, and the six controller boxes in each axis --- is three bytes:
an on/off and phase byte, a \textsc{function select} byte, and a depth. The
\textsc{function select} byte chooses what the controller moves, from this list.

\begin{center}
{\small
\begin{tabular}{llll}
\toprule
\bytes{00} & \textsc{off}                     & \bytes{1A} & filter LFO3 speed \\
\bytes{01} & pitch bend                       & \bytes{1B} & filter LFO4 speed \\
\bytes{02} & sustain level                    & \bytes{1C} & fitting$^\ast$ \\
\bytes{03} & filter cutoff frequency          & \bytes{1D} & position \\
\bytes{04} & pitch LFO1 depth                 & \bytes{1E} & position depth \\
\bytes{05} & pitch LFO2 depth                 & \bytes{1F} & (reserve) \\
\bytes{06} & pitch LFO3 depth                 & \bytes{20} & position movement width \\
\bytes{07} & pitch LFO4 depth                 & \bytes{21} & position movement speed \\
\bytes{08} & amp LFO1 depth                   & \bytes{22} & (reserve) \\
\bytes{09} & amp LFO2 depth                   & \bytes{23} & muting$^\ast$ \\
\bytes{0A} & amp LFO3 depth                   & \bytes{24} & resonator key shift$^\ast$ \\
\bytes{0B} & amp LFO4 depth                   & \bytes{25} & sub gain \\
\bytes{0C} & filter LFO1 depth                & \bytes{26} & delay$^\ast$ \\
\bytes{0D} & filter LFO2 depth                & \bytes{27} & panning$^\ast$ \\
\bytes{0E} & filter LFO3 depth                & \bytes{28} & effect1 send$^\ast$ \\
\bytes{0F} & filter LFO4 depth                & \bytes{29} & reverb send$^\ast$ \\
\bytes{10} & pitch LFO1 speed                 & \bytes{2A} & level$^\ast$ \\
\bytes{11} & pitch LFO2 speed                 & \bytes{2B} & amp env.\ attack time$^\ast$ \\
\bytes{12} & pitch LFO3 speed                 & \bytes{2C} & amp env.\ decay time$^\ast$ \\
\bytes{13} & pitch LFO4 speed                 & \bytes{2D} & amp env.\ release time$^\ast$ \\
\bytes{14} & amp LFO1 speed                   & \bytes{2E} & filter resonance$^\ast$ \\
\bytes{15} & amp LFO2 speed                   & \bytes{2F} & effect1 dynamic control \\
\bytes{16} & amp LFO3 speed                   & \bytes{30} & effect2 dynamic control \\
\bytes{17} & amp LFO4 speed                   & \bytes{31} & reverb dynamic control \\
\bytes{18} & filter LFO1 speed                & \bytes{32}--\bytes{7F} & (reserve) \\
\bytes{19} & filter LFO2 speed                &            & \\
\bottomrule
\end{tabular}
}
\end{center}

$^\ast$ marked functions cannot be driven by after touch.

\section{Filter values}\label{sec:filtervals}
A tone's filter has four \textsc{value} bytes whose meaning depends on the filter mode
selected for that tone.

\begin{center}
{\small
\begin{tabular}{lllll}
\toprule
mode & \textsc{value1} & \textsc{value2} & \textsc{value3} & \textsc{value4} \\
\midrule
\textsc{lpf12} & cutoff frequency & resonance & EQ frequency & EQ gain, EQ range \\
\textsc{hpf12} & cutoff frequency & resonance & EQ frequency & EQ gain, EQ range \\
\textsc{lpf24} & cutoff frequency & resonance & --- & --- \\
\textsc{hpf24} & cutoff frequency & resonance & --- & --- \\
\textsc{bpf}   & low cutoff frequency & low resonance & high cutoff frequency & high resonance \\
\bottomrule
\end{tabular}
}
\end{center}

Cutoff frequencies are \bytes{00}--\bytes{7F}. Resonance is \bytes{00}--\bytes{05}: in the
12\,dB modes and \textsc{bpf} a step is 3\,dB, covering $-6$ to $9$\,dB; in the 24\,dB
modes a step is 6\,dB, covering $-12$ to $18$\,dB. \textsc{eq gain} is
\bytes{00}--\bytes{0E}, where \bytes{00} is silence and \bytes{01}--\bytes{0E} span $-6$ to
$+6$\,dB, and the top bit of that same byte is \textsc{eq range}: \bytes{0} low,
\bytes{1} high.

\section{Interaction connect}
Modeling parameter \bytes{228} sets how a sound's four resonators are connected --- none of
them, the first pair, the second pair, both pairs, or all four together. The guide gives
the resulting connection state of each of the four blocks for every setting; it is not
reproduced here, and the parameter's own range, \bytes{00}--\bytes{03}, is in
chapter~\ref{ch:sound}.

\chapter{Universal System Exclusive}

\section{Recognised messages}
The instrument honours two universal messages, and transmits both.

\begin{center}
\begin{tabular}{ll}
\toprule
message & meaning \\
\midrule
\bytes{F0 7E 7F 09 01 F7} & General MIDI System On \\
\bytes{F0 7E 7F 09 02 F7} & General MIDI System Off \\
\bottomrule
\end{tabular}
\end{center}

Turning General MIDI on restricts the sounds and operations available and may rearrange
the percussion map. General MIDI System Off has no effect if General MIDI is already off.

Both messages are also written into Standard MIDI Files produced by the instrument, as a
five-byte System Exclusive event at delta time zero.

\section{Messages not recognised}
A controller expecting any of the following will receive no response:

\begin{center}
\begin{tabular}{ll}
\toprule
message & \\
\midrule
\bytes{F0 7F 7F 04 01 \textit{ll} \textit{mm} F7} & Master Volume \\
\bytes{F0 7E \textit{dd} 06 01 F7} & Identity Request \\
\bytes{F0 7E 7F 09 03 F7} & General MIDI 2 System On \\
\bottomrule
\end{tabular}
\end{center}

The instrument does not implement Identity Reply, so it cannot be auto-detected by a
controller that relies on it. The entire Universal Real Time identifier \bytes{7F} is
unrecognised.

\chapter{Settings that affect System Exclusive}\label{ch:settings}

\section{The EXCLUSIVE filter}
The MIDI menu's \textsc{input \& output filter} page has eight rows, each a plain
\textsc{on}/\textsc{off} switch for one kind of MIDI message. The last row is
\textsc{exclusive}. There is one such switch and it governs both directions: the same
setting is consulted when the instrument sends and when it receives. There is no separate
input filter, no per-socket filter and no third state, and the page is the same on both
models.

\begin{caution}{EXCLUSIVE governs the tempo message and nothing else}
Despite its name, this switch does not turn System Exclusive on and off. Its entire effect
on the messages in this reference is the family \bytes{25} tempo message: with
\textsc{exclusive} set to \textsc{off} the instrument neither sends one nor acts on one
received, silently in both directions.

\textbf{A bulk dump transfers normally with \textsc{exclusive} set to \textsc{off}.}
Software that cannot complete a transfer should not look here for the cause.
\end{caution}

Everything below is unaffected by the setting:

\begin{center}
\begin{tabular}{ll}
\toprule
function & affected by \textsc{exclusive}? \\
\midrule
bulk dump, sent or received & no \\
dump requests & no \\
the opening enquiry and its reply & no \\
acknowledgements and the session control messages & no \\
parameter messages, and the replies and refusals to them & no \\
General MIDI System On / Off, sent or received & no \\
error reporting & no \\
\midrule
tempo message (family \bytes{25}) & \textbf{yes} \\
\bottomrule
\end{tabular}
\end{center}

\section{Which sockets carry what}
The instrument has two sets of MIDI sockets. Only bulk-dump traffic is restricted to the
first.

\begin{center}
\begin{tabular}{p{88mm}l}
\toprule
message & accepted on \\
\midrule
bulk-dump data, continuation frames, handshake, acknowledgements & MIDI 1 only \\
dump requests & MIDI 1 only \\
parameter messages (families \bytes{2B}, \bytes{2C}) & MIDI 1 and MIDI 2 \\
tempo message (family \bytes{25}) & MIDI 1 and MIDI 2 \\
General MIDI System On / Off & MIDI 1 and MIDI 2 \\
\bottomrule
\end{tabular}
\end{center}

MIDI channels are written with the socket set first: \texttt{1-01} to \texttt{1-16} and
\texttt{2-01} to \texttt{2-16}. No channel or device-number setting affects System
Exclusive addressing --- these messages carry no channel.

\section{Memory protect}
Two write-protect switches are consulted before an incoming message is acted on. A message
they block never reaches the transfer machinery, and the instrument shows its own
write-protection prompt rather than a System Exclusive error. A transfer that stops with
that prompt on screen has not failed as a transfer.

\chapter{Errors}

Transfer failures are reported on the \textsc{sysex bulk dump} screen. A System Exclusive
message arriving while any other screen is displayed produces no visible error.

\section{The messages}
\begin{center}
\begin{tabular}{ll}
\toprule
screen & meaning \\
\midrule
\textsc{completed!} & the transfer finished \\
\textsc{error 40!} & the message is addressed to a different product \\
\textsc{error 41!} & reception failed \\
\textsc{error 42!} & transmission failed \\
\textsc{error 21!} & the destination is full \\
\bottomrule
\end{tabular}
\end{center}

\begin{caution}{Not every failure is a 40, 41 or 42}
When incoming data is larger than the destination area, the instrument reports
\textsc{error 21! Memory full}. Software that recognises only the three documented
System Exclusive errors will mis-read a truncated restore as something else.
\end{caution}

\section{ERROR 40 --- wrong identification}
Raised when the message family byte is accepted but the bytes that follow it do not match
the product. For a message addressed to the SX-WSA1R this means the model triple was not
\bytes{04}, then \bytes{00} or \bytes{01}, then \bytes{11}.

\section{ERROR 41 --- reception}
The general reception failure. It covers:

\begin{itemize}
\item a message inside a session that does not begin \bytes{F0 50};
\item a status byte other than \bytes{F7} arriving inside an open message;
\item a framing, parity or overrun error on the MIDI input, or loss of active sensing;
\item no reply within the timeout;
\item a message longer than 254 bytes;
\item an odd number of payload bytes;
\item a continuation flag that is neither \bytes{00} nor \bytes{01};
\item a checksum mismatch.
\end{itemize}

\section{ERROR 42 --- transmission}
Raised when the instrument cannot complete an outgoing transfer internally, roughly a
second after the attempt stalls. In this version it has a single cause and does not
indicate a fault on the MIDI cable or at the receiving device.

\section{Interference from other traffic}
While a bulk-dump session is open, unrelated System Exclusive traffic arriving on the MIDI
input can set the error state, and will then be reported when the session ends. Keep other
sources off the MIDI~1 input for the duration of a transfer.

\chapter{SX-WSA1 and SX-WSA1R}\label{ch:models}

The SX-WSA1 keyboard and the SX-WSA1R rack module run the same program and speak the same
System Exclusive protocol. Two features are nonetheless available on the keyboard and not
on the rack, in both directions.

\begin{center}
\begin{tabular}{lcc}
\toprule
 & SX-WSA1 & SX-WSA1R \\
 & (keyboard) & (rack) \\
\midrule
message framing, checksum, handshake & yes & yes \\
\textsc{sound} bulk dump & yes & yes \\
\textsc{combination} bulk dump & yes & yes \\
\textsc{system, part \& midi} bulk dump & yes & yes \\
dump requests & yes & yes \\
General MIDI On / Off & yes & yes \\
\midrule
\textsc{sequencer} bulk dump & yes & \textbf{no} \\
tempo message (family \bytes{25}) & yes & \textbf{no} \\
\bottomrule
\end{tabular}
\end{center}

\begin{caution}{On the rack, a SEQUENCER dump sends nothing and says nothing}
The \textsc{sequencer} row is present on the rack's \textsc{sysex bulk dump} menu, and
selecting it and pressing \textsc{send} appears to work. Nothing is transmitted --- not
even the end-of-category message that closes every other transfer. A librarian waiting for
data will simply time out.

The same applies to the \textsc{sequencer} portion of \textsc{total keyboard}.
\end{caution}

The tempo message is refused in both directions on the rack: it will neither act on one
received nor send one.

\begin{note}{The model triple does not tell you which you are talking to}
Both models transmit the same model triple, so it cannot be used to distinguish them. The
only reliable test over MIDI is behavioural --- request a \textsc{sequencer} dump and see
whether anything arrives.
\end{note}

\chapter{Relationship to other Technics instruments}

The message engine described in this reference is not unique to the SX-WSA1R. At least two
other Technics instruments implement the same protocol: the SX-KN5000 and the SX-KN1500.

\section{What is shared}
All three use the \bytes{50} identifier, the same message families, the same short control
messages, the same payload encoding and the same checksum rule. The fixed messages are
byte-for-byte identical:

\begin{center}
\bytes{F0 50 23 7E F7}\quad
\bytes{F0 50 24 7E F7}\quad
\bytes{F0 50 27 7E F7}\quad
\bytes{F0 50 28 7E F7}\quad
\bytes{F0 50 29 7E F7}\quad
\bytes{F0 50 2A 7E F7}
\end{center}

The error mapping is also nearly common: of the thirty-four status values, thirty-two
produce the same report on the SX-WSA1R and the SX-KN5000.

\section{What distinguishes the products}
\textsc{pc}, \textsc{md} and \textsc{ver} together are the discriminator. A message whose
three bytes do not match the receiving instrument is refused with \textsc{error 40}.

\begin{center}
\begin{tabular}{ll}
\toprule
instrument & \textsc{pc} \textsc{md} \textsc{ver} \\
\midrule
SX-WSA1 / SX-WSA1R & \bytes{04 00 11}, \bytes{04 01 11} \\
SX-KN5000 & \bytes{01 28 12} \\
SX-KN1500 & \bytes{01 24 11}, \bytes{01 25 11}, \bytes{01 26 11} \\
\bottomrule
\end{tabular}
\end{center}

The split is legible once the fields are named: \textsc{pc}~=~\bytes{04} is the
\textsc{wsa} product category and \bytes{01} the keyboards, \textsc{md} separates models
within a category, and \textsc{ver} is the exclusive specification's own version --- 2.1
for the WSA1 and the KN1500, 2.2 for the KN5000.

A librarian that supports more than one of these instruments needs only to vary those bytes
and the dump region addresses; the framing, handshake, encoding and checksum are common
code.

\section{Dump regions differ}
The dump region addresses are per-product. The SX-KN5000 and SX-KN1500 share eleven of
their thirteen region and length pairs with each other, but the SX-WSA1R shares none of
its eight with either --- its categories are its own.

\begin{caution}{The SX-WSA1R does not accept the GS-style messages its siblings do}
The SX-KN5000 and SX-KN1500 additionally accept a small set of messages in the
\bytes{42 12} parameter form, including the sequence \bytes{42 12 40 00 7F 00 41} used
elsewhere in the industry as a reset. \textbf{The SX-WSA1R accepts none of them.} Software
that resets a KN-series instrument that way and assumes the SX-WSA1R will follow will find
it does not respond.
\end{caution}

\chapter{Worked examples}\label{ch:examples}

Three things in this format are easy to get subtly wrong: the seven-bit address encoding,
the high-nibble-first payload, and which bytes the checksum covers. The examples below are
complete messages, and every byte in them is computed rather than transcribed.

%% GENERATED by notes/sysex-probes/sysex_examples_tex.py -- do not edit

\section{A parameter, written and read back}
The example is part~0's parameter \bytes{03}, a one-byte value, set to 100.

\begin{center}
\begin{tabular}{ll}
\toprule
& message \\
\midrule
write & \bytes{F0 50 2C 04 00 11 00 20 03 00 00 01 06 04 00 41 F7} \\
request & \bytes{F0 50 2B 04 00 11 00 20 03 00 00 01 00 4C F7} \\
reply & \bytes{F0 50 2C 04 00 11 00 20 03 00 00 01 06 04 00 41 F7} \\
part 7 & \bytes{F0 50 2C 04 00 11 00 27 03 00 00 01 06 04 00 3A F7} \\
three-byte & \bytes{F0 50 2C 04 00 11 00 20 00 00 00 03 01 02 00 00 04 00 00 45 F7} \\
\bottomrule
\end{tabular}
\end{center}

The value 100 is \bytes{64}, which the payload carries as the two bytes
\bytes{06 04}. The three bytes before it are the count, and the \bytes{00} after
it is the continuation flag. The checksum closes the message.

\section{A frame from a real instrument}
The bytes below open the first data message of a dump recorded from an
SX-WSA1R, a transfer of three categories in one session:

\begin{center}
\bytes{F0 50 2D 04 01 11 40 00 00 00 00 20 \ldots}
\end{center}

Reading it by the rules of chapter~4: family \bytes{2D}, model triple
\bytes{04 01 11} --- the rack's --- destination \bytes{40 00 00} and length \bytes{00 00 20},
which together name block~1 of \textsc{system, part \& midi}. The length decodes to 32 source
bytes and the frame carries 32 of them.

The payload decodes to \bytes{5A 5A 01 00 57 41 30 00}\ldots, the signature chapter~5 gives for
that block.

That block is the one chapter~\ref{ch:params} maps: a dump of this category can be
read against table~\ref{tbl:partparams} directly, which is the cheapest way to
confirm a parameter address without an instrument to hand.


\chapter{Scope of this edition}\label{ch:limits}

This reference covers message framing, the bulk-dump transfer protocol, the parameter
address space, the universal messages, error reporting and the differences between the two
models. The following are not covered.

\begin{description}
\item[What the parameters do.] Chapter~\ref{ch:params} enumerates all 108 addresses of
  the \bytes{00} area with their name, length, bit mask, accepted values and position in a
  dump. It does not explain what each setting does to the sound; the user manuals do that.
\item[The effect parameters.] The twenty \textsc{value} bytes under each of
  \textsc{effect1}, \textsc{effect2} and \textsc{reverb} mean different things for each
  effect type. The published guide gives them, effect by effect, and this edition does not
  reproduce them.
\item[Block contents.] Chapter~\ref{ch:blocks} gives each block's signature and where its
  names are, but not how a stored sound, combination or sequence is arranged inside one.
  Chapter~\ref{ch:sound} lays out the sound being \emph{edited}, which is a different
  thing from the sound memory a bulk dump carries.
\item[The SEQUENCER category, against hardware.] Its protocol is described, but the
  capture used to check everything else came from a rack, which cannot send it. See
  chapter~\ref{ch:overview}.
\item[Operating system version 1.] Everything here describes version 2.0.
\end{description}


\end{document}
